\documentclass[11pt,a4paper]{article}
\usepackage[margin=1in]{geometry}
\usepackage[utf8]{inputenc}
\usepackage[english]{babel}
\usepackage{amsmath,amssymb}
\usepackage{microtype}
\usepackage{hyperref}

\setlength{\parindent}{0pt}
\setlength{\parskip}{0.75\baselineskip}
\sloppy

\title{Horizon as a Phase Recording Surface in Field 01\\\large A Black Hole Without a Point Singularity and Without Literal Evaporation}
\author{}
\date{}

\newcommand{\field}{\mathcal{F}_{01}}
\newcommand{\phase}{\varphi}
\newcommand{\normal}{\mathcal{N}}
\newcommand{\memory}{\mathcal{M}}
\newcommand{\bulk}{\mathrm{bulk}}
\newcommand{\boundary}{\mathrm{boundary}}

\begin{document}
\maketitle

\begin{center}
\textit{Working version v0.1 --- English research draft}
\end{center}

\begin{abstract}
This paper formulates the horizon block of Field 01. A black hole is interpreted not as an object containing a physical point singularity, but as a limiting regime of the field in which a volumetric phase structure is transformed into a surface record. In this language the horizon is not an entrance to infinite density, but a boundary at which the radial scalar/normal profile is suppressed, volumetric unfolding of the dynamics ceases, and phase memory is retained. The Bekenstein--Hawking entropy is interpreted as an indication of the surface character of the record. At the same time, Field 01 differs from a literal evaporation picture of black holes: the horizon does not lose its record as an ordinary thermal emitter, but accumulates phase information. This remains a hypothesis and requires a strict comparison with Hawking's semiclassical theory, quantum field theory in curved spacetime, and holographic approaches.
\end{abstract}

\textbf{Keywords:} Field 01; black hole; horizon; phase record; memory; Bekenstein--Hawking entropy; holography; singularity; Hawking radiation; quantum field theory.

\section{How to Read This Paper}

This paper continues the work on the elementary particle as a closed wave of Field 01. In the first paper, a local particle was described as a closed phase node with a retained scalar/normal profile, mass, and memory. Here we consider the limiting case: what happens when such a local structure reaches a black-hole horizon.

The text should be read on three levels:
\begin{enumerate}
\item \textbf{Established physics}: in general relativity black holes possess horizons; in black-hole thermodynamics entropy is proportional to horizon area; in Hawking's semiclassical theory the horizon is associated with thermal radiation.
\item \textbf{Interpretation proposed by Field 01}: the horizon is understood as a recording surface on which volumetric phase structure is translated into a surface code.
\item \textbf{Open hypothesis}: Field 01 does not adopt literal black-hole evaporation as loss of record through a thermal flux; instead, it assumes accumulation of phase information on the horizon.
\end{enumerate}

The paper therefore does not claim that Hawking's standard calculations have been refuted. It formulates an alternative language that must be compared separately with semiclassical horizon physics.

\section{What Is Accepted and What Is Reinterpreted}

To avoid misreading, we fix the boundary between standard physics and Field 01. The model does not deny the role of the horizon, entropy, or reduced descriptions. It reinterprets the physical meaning of these structures.

\begin{center}
\begin{tabular}{p{0.43\textwidth}p{0.43\textwidth}}
\textbf{Standard description} & \textbf{Field 01 interpretation} \\
the horizon limits causal access & the horizon is a surface of phase recording \\
entropy is proportional to area & area measures the capacity of the surface archive \\
an external observer sees a reduced state & only part of the phase record is externally accessible \\
Hawking radiation is thermal & thermality is read as an effect of limited access \\
evaporation describes mass loss through a flux & literal loss of record is not assumed; the record accumulates on the horizon \\
a singularity appears in the classical inward continuation & the point center is replaced by limiting surface fixation.
\end{tabular}
\end{center}

This table is not a proof of the model. It is a map of differences: what Field 01 accepts as an important physical fact and what it translates into its own language.

\section{Main Idea}

In Field 01 a black hole is not understood as a hole with a physical point of infinite density inside. It is understood as a limiting regime of the field:
\begin{equation}
\text{BH} \neq \text{point singularity inside},
\qquad
\text{BH} \sim \text{surface of limiting fixation}.
\end{equation}

A key formulation of the model is:
\begin{equation}
\text{zero is not inside; zero is the boundary.}
\end{equation}

This means that the limiting regime is not placed at the center as a small region of infinite density. Rather, the very concept of internal volume ceases to be primary. The physically significant structure becomes the horizon as the surface on which the phase structure is fixed.

\section{Local Particle Memory and Global Horizon Record}

In Field 01 an elementary particle is a local node of memory. It retains a phase structure through closed circulation and a retained scalar/normal profile. Schematically,
\begin{equation}
\text{particle} = \text{local memory} + \normal.
\end{equation}

The black-hole horizon is another regime of memory. It does not retain the local scalar/normal profile of an individual particle. Instead, it translates many local phase structures into a global surface record:
\begin{equation}
\memory_{\mathrm{particle}} \longrightarrow \memory_{\mathrm{horizon}}.
\end{equation}

In terms of volume and boundary, this transition can be written as
\begin{equation}
\memory_{\bulk} \longrightarrow \memory_{\boundary}.
\end{equation}

This expression is not yet a strict mapping operator. It fixes the principle that must be formalized: phase information associated with a volumetric structure should be represented on the horizon surface.

\section{Relation to Bekenstein--Hawking Entropy}

The entropy of a black hole is given by
\begin{equation}
S_{\mathrm{BH}} = \frac{k_B A}{4\ell_P^2},
\end{equation}
where $A$ is the horizon area, $\ell_P$ is the Planck length, and $k_B$ is Boltzmann's constant.

For standard physics this formula expresses a deep relation between the horizon, thermodynamics, quantum theory, and gravity. For Field 01 the important point is that the entropy scales with area rather than with volume. In the language of the model this is read as an indication of the surface character of the record:
\begin{equation}
\text{number of accessible records} \sim A.
\end{equation}

This does not derive the Bekenstein--Hawking formula. It provides an interpretive reading: if the horizon is a phase recording surface, then the area law becomes natural as a statement about the capacity of that surface.

\section{Why a Point Singularity Is Not Required}

The point singularity appears in the classical continuation of the general-relativistic solution. Field 01 does not treat this continuation as a literal physical object. The model proposes another endpoint of collapse: not concentration into a point, but translation into a boundary record.

The difference can be written schematically:
\begin{equation}
\text{collapse to point} \quad \longrightarrow \quad \text{collapse of the local scalar/normal profile}.
\end{equation}

If the local scalar/normal profile is suppressed,
\begin{equation}
\normal \rightarrow 0,
\end{equation}
then the local volumetric node can no longer remain an ordinary particle-like structure. Its phase memory must either disappear or be represented differently. Field 01 chooses the second option:
\begin{equation}
\normal \rightarrow 0,
\qquad
\memory_{\bulk} \rightarrow \memory_{\boundary}.
\end{equation}

Thus the singular point is replaced by a limiting surface process. The model does not claim that this is already a complete quantum-gravity solution. It proposes a consistent interpretation to be formalized.

\section{What Happens to a Closed Wave at the Horizon}

A closed wave has phase circulation and a local scalar/normal profile. Near the horizon the radial direction becomes limiting for the external description. In the language of Field 01, the local scalar/normal profile is suppressed and tangential or surface degrees of freedom become dominant.

Schematically,
\begin{equation}
\text{closed wave with } \normal \neq 0
\quad \longrightarrow \quad
\text{surface phase record with } \normal \to 0.
\end{equation}

This transition should not be imagined as a mechanical tearing of the particle. The particle is not a solid body. It is a phase structure. Therefore the limiting process is better described as a change of regime:
\begin{equation}
\text{local phase node} \longrightarrow \text{boundary phase code}.
\end{equation}

The local object-like description disappears, but the phase pattern is not assumed to be destroyed.

\section{Difference from Literal Evaporation}

In semiclassical physics Hawking radiation is derived as a thermal flux associated with quantum fields on a black-hole background. Field 01 does not deny that an external observer may describe a thermal reduced state. The difference concerns the interpretation.

In Field 01 thermality is read as a consequence of restricted access to the full phase record. The external observer sees only a reduced description:
\begin{equation}
\rho_{\mathrm{external}} = \mathrm{Tr}_{\mathrm{hidden}}\, \rho_{\mathrm{full}}.
\end{equation}

The thermal character of this reduced state does not necessarily mean that the horizon literally loses the underlying phase record. The model therefore separates two statements:
\begin{enumerate}
\item an external thermal description may exist;
\item the fundamental phase record is not assumed to be destroyed by that description.
\end{enumerate}

This is the cautious formulation of the difference. The model still must explain how the standard Hawking result emerges and where, if anywhere, it predicts a different behavior.

\section{Horizon as a Phase Archive}

The horizon can be understood as an archive of phase relations. The word archive does not mean a material storage device placed on a surface. It means that the accessible physical description becomes boundary-based.

A schematic notation is
\begin{equation}
\mathcal{A}_{\mathrm{horizon}} = \left\{ \memory_i \right\}_{i=1}^{N(A)},
\end{equation}
where $N(A)$ is the number of distinguishable records associated with a horizon of area $A$.

The exact relation between $N(A)$ and $S_{\mathrm{BH}}$ is an open problem. One possible direction is
\begin{equation}
S_{\mathrm{BH}} \sim k_B \log N(A).
\end{equation}

This would connect the Field 01 language of phase records with the statistical interpretation of entropy.

\section{The Last Closed Wave and the Cycle}

The book introduces the limiting image of the last elementary particle, or the last closed wave. This is not a claim about a specific observable particle. It is a cosmological image of the last remaining distinction.

As long as at least one closed node exists, the following remain:
\begin{itemize}
\item the distinction between regimes;
\item the direction $1$;
\item a minimal phase delay;
\item the possibility of counting changes, that is, time.
\end{itemize}

When the last closed wave passes into the recording regime, what disappears is not an object in the ordinary sense, but the last retained phase delay:
\begin{equation}
\text{last particle} \longrightarrow \text{last horizon record}.
\end{equation}

After that, the external regime no longer retains distinctions. In the book this is described as the limiting transition $0$--$0$: the coincidence of two fixation regimes, from which a new minimal impulse of distinction may arise.

This section remains the most speculative part of the paper. Its purpose is to preserve the connection with the cosmological part of the book without mixing it with the more cautious physical part of the horizon article.

\section{What Must Be Formalized}

For the horizon block of Field 01 to become a physical theory, one must define:
\begin{enumerate}
\item the horizon state space $\mathcal{H}_{\mathrm{horizon}}$;
\item the map $\memory_{\bulk}\to\memory_{\boundary}$;
\item a conservation law for phase information;
\item the relation between horizon area and the number of distinguishable phase records;
\item the relation of the model to Hawking's calculation;
\item observational or theoretical differences from the standard evaporation picture;
\item the connection with the holographic principle and the black-hole information problem.
\end{enumerate}

Until these steps are completed, Field 01 should remain cautious: it offers a language and a hypothesis, not a completed proof.

\section{Conclusion}

Field 01 proposes to interpret a black hole as a limiting regime of the field in which local phase structures lose their local scalar/normal profile and become surface records. In this description the horizon is not an entrance to a point singularity, but an archive of phase information.

The main difference from literal evaporation is the following: the horizon does not destroy or lose the record as an ordinary thermal emitter. It limits external access to information and translates volumetric structure into a surface code. This is consistent with the area motif in the Bekenstein--Hawking formula, but it requires a separate strict formalization and comparison with semiclassical calculations.

The next step is to build a mathematical map between volumetric memory and surface record, and to clarify how horizon thermality arises as a reduced description rather than as literal disappearance of the phase archive.

\end{document}